\documentclass[11pt]{article}
\usepackage[margin=2.5cm]{geometry}
\usepackage{amsmath,amssymb,bm}
\usepackage{booktabs}
\usepackage{graphicx}
\usepackage{hyperref}
\usepackage{xcolor}
\hypersetup{colorlinks=true, linkcolor=blue!60!black, citecolor=blue!60!black,
            urlcolor=blue!60!black}

\newcommand{\eps}{\varepsilon}
\newcommand{\lm}{\lambda_{\ell m}}
\newcommand{\cth}{\cos\theta}
\newcommand{\sth}{\sin\theta}

\title{\textbf{Almond}: a GPU spherical-harmonic transform for HEALPix\\
{\large Mathematical techniques and technical design}}
\author{PowerSpec / SiMaster project}
\date{July 2026 (prototype v0.1: spin-0 synthesis)}

\begin{document}
\maketitle

\begin{abstract}
\noindent
\textbf{Almond} is a CUDA implementation of the spherical-harmonic transform
(SHT) on HEALPix grids, intended to replace the CPU library \texttt{ducc0}
as the SHT backend of the SiMaster quadratic maximum-likelihood (QML) power
spectrum estimator.  It implements, on the GPU, the same
exact algorithm as \texttt{ducc0} (the libsharp2 lineage): double precision
throughout, healpy conventions, matrix-free, with the every-other-$\ell$
Legendre recursion, $2^{\pm800}$ dynamic-range scaling, and per-ring FFT-bin
aliasing.  \texttt{ducc0} is the correctness reference; the prototype
achieves $\lesssim 10^{-11}$ maximum relative deviation from it.  This
document derives the algorithm, records the design decisions and their
GPU-specific rationale, and reports benchmark results on NERSC Perlmutter
A100 nodes.
\end{abstract}

\tableofcontents

%==============================================================================
\section{Problem statement and conventions}

A real spin-0 field sampled on a HEALPix RING grid of resolution
$N_\mathrm{side}$ is synthesised from harmonic coefficients as
\begin{equation}
f(\theta,\varphi) \;=\; \sum_{\ell=0}^{\ell_{\max}} \sum_{m=-\ell}^{\ell}
   a_{\ell m}\, Y_{\ell m}(\theta,\varphi),
\qquad
Y_{\ell m}(\theta,\varphi) = \lm(\theta)\, e^{im\varphi},
\end{equation}
where $\lm$ are the fully normalised associated Legendre functions
($\int |Y_{\ell m}|^2 d\Omega = 1$).  healpy stores only $m \ge 0$; reality
implies $a_{\ell,-m} = (-1)^m a_{\ell m}^{*}$ and
$\lambda_{\ell,-m} = (-1)^m \lm$.  The triangular healpy layout indexes
$a_{\ell m}$ at $\mathrm{mstart}(m) + \ell$ with
$\mathrm{mstart}(m) = m(2\ell_{\max}+1-m)/2$.  Following
\texttt{ducc0}/healpy, the imaginary part of $a_{\ell 0}$ is ignored.

HEALPix RING geometry places pixels on $4N_\mathrm{side}-1$ iso-latitude
rings; ring $i$ ($1 \le i \le 4N_\mathrm{side}-1$, north to south) has
\begin{equation}
n_\varphi(i) = \begin{cases} 4i' & \text{polar caps } (i' < N_\mathrm{side})\\
4N_\mathrm{side} & \text{equatorial belt} \end{cases}
\qquad
\varphi_j = \varphi_0(i) + \frac{2\pi j}{n_\varphi},
\end{equation}
with $i' = \min(i, 4N_\mathrm{side}-i)$, and $\cth$ an exact rational:
$1 - i'^2/(3N_\mathrm{side}^2)$ in the caps,
$4/3 - 2i/(3N_\mathrm{side})$ in the belt.  We compute $\cth$, $\sth$
directly from these rationals (never through $\theta=\arccos z$), and
$\varphi_0 = \pi q/d$ with small integers $(q,d)$ --- both facts are
exploited for accuracy (\S\ref{sec:fold}).

The transform factorises over rings:
\begin{equation}
F_m(\theta) \;=\; \sum_{\ell=m}^{\ell_{\max}} a_{\ell m}\, \lm(\theta),
\qquad
f(\theta,\varphi_j) \;=\; \sum_{m=-\ell_{\max}}^{\ell_{\max}}
   F_m(\theta)\, e^{im\varphi_j},
\quad F_{-m} = F_m^{*}.
\label{eq:factorize}
\end{equation}
Stage 1 (\emph{Legendre}, $O(N_\mathrm{ring}\,\ell_{\max}^2)$ work) computes
$F_m$ on every ring; stage 2 (\emph{FFT}, $O(N_\mathrm{pix}\log
N_\mathrm{side})$) evaluates the azimuthal sums.  Stage 1 dominates at all
resolutions of interest.

%==============================================================================
\section{Stage 1: the every-other-$\ell$ Legendre recursion}
\label{sec:legendre}

\subsection{Three-term recurrence and ring pairing}

The normalised Legendre functions obey
\begin{equation}
\cth\, \lm = \eps_{\ell+1}\lambda_{\ell+1,m} + \eps_{\ell}\lambda_{\ell-1,m},
\qquad
\eps_\ell = \sqrt{\frac{\ell^2 - m^2}{4\ell^2 - 1}},
\label{eq:threeterm}
\end{equation}
seeded by $\lambda_{mm}(\theta) = c_m \sth^m$,
$c_m = (-1)^m\sqrt{(2m+1)!!/(4\pi\,(2m)!!)}$, and
$\lambda_{m-1,m}=0$.  The parity relation
$\lm(\pi-\theta) = (-1)^{\ell+m}\lm(\theta)$ makes it natural to process
rings in north/south pairs sharing $|\cth|$: splitting the sum in
\eqref{eq:factorize} into even and odd $\ell-m$,
\begin{equation}
S_e = \!\!\sum_{\ell-m\ \mathrm{even}}\!\! a_{\ell m}\lm, \qquad
S_o = \!\!\sum_{\ell-m\ \mathrm{odd}}\!\! a_{\ell m}\lm, \qquad
F_m^{\rm north} = S_e + S_o,\quad F_m^{\rm south} = S_e - S_o .
\label{eq:evenodd}
\end{equation}
This halves the Legendre work.  The HEALPix equator ring is self-paired
($\cth=0$).

\subsection{Folding the odd $\ell$ into the even stream}
\label{sec:fold-l}

\texttt{ducc0}'s key trick (inherited here) is to visit only every other
$\ell$ in the $\theta$ recursion.  Write $n_{i} = m + 2i + 1$ for the odd
offsets, $x = \cos^2\theta$.  Applying \eqref{eq:threeterm} twice gives a
pure $\Delta\ell = 2$ recurrence,
\begin{equation}
\lambda_{n+2} \;=\;
\frac{(x - \eps_n^2 - \eps_{n+1}^2)\,\lambda_{n}
      \;-\; \eps_{n-1}\eps_{n}\,\lambda_{n-2}}{\eps_{n+1}\,\eps_{n+2}} .
\label{eq:delta2}
\end{equation}
Introduce scaled variables $\mu_i$ and per-$m$ constants $\alpha_i$:
\begin{equation}
\mu_i \;=\; \frac{\lambda_{n_i}}{\alpha_i \cth},
\qquad
\alpha_0 = \frac{1}{\eps_{m+1}},\qquad
\alpha_{i+1} = \frac{\sigma_i}{\eps_{n_i+1}\,\eps_{n_i+2}\,\alpha_i},
\qquad \sigma_i = (-1)^{i} .
\label{eq:alpha}
\end{equation}
Then \eqref{eq:delta2} becomes a recurrence with \emph{unit} trailing
coefficient --- one FMA for the polynomial, one for the update:
\begin{equation}
\boxed{\;\mu_{i+1} = (a_i\,x + b_i)\,\mu_i + \mu_{i-1}\;},
\qquad
a_i = \sigma_i\,\alpha_i^2,\qquad
b_i = -a_i\,(\eps_{n_i}^2 + \eps_{n_i+1}^2),
\label{eq:murec}
\end{equation}
with exact initial conditions $\mu_{-1} = 0$ and
$\mu_0 = \lambda_{mm}$ (because $\lambda_{m+1,m} = \cth\,\lambda_{mm}/
\eps_{m+1}$).  The $\alpha_i$ carry the sign pattern
$(+,+,-,-,+,+,\dots)$ of period four and $|\alpha_i| = \sqrt{|a_i|}$, which
the implementation uses to reconstruct $\alpha$ from the stored $a$
(\S\ref{sec:tables}).  Note that $\cth$ never appears in a denominator: at
the equator the recursion runs on finite $\mu_i$ even though
$\lambda_{n_i}(\pi/2) = 0$.

Both partial sums come out of the \emph{same} $\mu$ sequence.  Pre-fold the
coefficients once per transform ($O(\ell_{\max}^2)$ total, embarrassingly
parallel):
\begin{equation}
A_i = \alpha_i\left(\eps_{n_i}\, a_{n_i-1,m} + \eps_{n_i+1}\, a_{n_i+1,m}\right),
\qquad
B_i = \alpha_i\, a_{n_i,m} .
\label{eq:prefold}
\end{equation}
Using \eqref{eq:threeterm} and reindexing (all boundary terms vanish since
$\lambda_{m-1,m}=0$),
\begin{equation}
\sum_i \mu_i A_i
 = \frac{1}{\cth}\sum_{\ell-m\,\mathrm{even}} a_{\ell m}
   \left(\eps_{\ell+1}\lambda_{\ell+1} + \eps_{\ell}\lambda_{\ell-1}\right)
 = S_e,
\qquad
\cth \sum_i \mu_i B_i = S_o .
\end{equation}
Hence, per ring pair and per $m$,
\begin{equation}
p_1 = \sum_i \mu_i A_i,\quad p_2 = \sum_i \mu_i B_i,
\qquad
F_m^{\rm north} = p_1 + \cth\, p_2,\quad
F_m^{\rm south} = p_1 - \cth\, p_2 .
\end{equation}
Cost: covering two $\ell$'s takes 2 real FMAs of recursion and 4 of
accumulation ($p_{1,2}$ complex) --- \emph{3 flops per $(\mathrm{ring\,pair},
m,\ell)$}, vs.\ 4--6 for the naive $\Delta\ell=1$ recursion, with half the
serial dependency chain.

\subsection{Dynamic range: $2^{\pm 800}$ scaling}
\label{sec:scaling}

$\lambda_{mm} \propto \sth^m$ underflows double precision catastrophically
near the poles ($\sth^{6000}$ at $N_\mathrm{side}=2048$ polar rings can be
$10^{-10^{4}}$).  Two devices (both from \texttt{ducc0}):

\begin{itemize}
\item \textbf{Scaled representation.}  A value is stored as
  $(v, s)$ meaning $v \cdot 2^{800 s}$ with $|v|$ kept below
  $2^{-60}$; the seed $\sth^m$ is computed by binary exponentiation in this
  representation (\texttt{mypow}, $O(\log m)$ multiplies --- accurate to
  $\sim\!10^{-15}$ where naive \texttt{pow} would lose $m\cdot$ulp).  While
  $s<0$ the true value is $< 2^{-860}$ and contributes exactly $0$ at double
  precision: the GPU thread advances the recursion \emph{without
  accumulating} and re-scales ($v\!\cdot\!2^{-800}$, $s{+}{+}$) whenever
  $|v|$ exceeds $2^{-60}$.  Once $s=0$ the thread switches permanently to
  the plain FMA loop (\S\ref{sec:kernel}); normalised $\lm$ are bounded by
  $\sim\!\ell^{1/4}$ so no overflow re-scaling is ever needed there.
\item \textbf{$m$ cutoff.}  On a ring with $\sth$ small, all $m >
  m_\mathrm{lim}(\theta) \approx \ell_{\max}\sth + \max(100,
  0.01\,\ell_{\max})$ give $F_m$ below double-precision relevance; those
  $(\mathrm{ring},m)$ pairs are skipped entirely and their $F_m$ set to $0$
  (identical rule to \texttt{ducc0}, so agreement with the reference is
  exact).  This removes $\sim 30\%$ of the Legendre work.
\end{itemize}

\subsection{GPU kernel design}
\label{sec:kernel}

CPU \texttt{ducc0} vectorises over rings within a SIMD register (AVX
lanes share one coefficient stream) and threads over $m$.  The GPU analogue
maps \textbf{one thread to two adjacent (ring pair, $m$) chains}:

\begin{itemize}
\item Grid: $y$-dimension $= m \in [0, m_{\max}]$; $x$-dimension covers
  the $2N_\mathrm{side}$ ring pairs in blocks of 128 threads.  At
  $N_\mathrm{side}=2048$ this is $\sim\!2\cdot10^5$ blocks / $2.5\cdot10^7$
  threads, each running $(\ell_{\max}-m)/2$ recursion steps: ample
  parallelism, and the strongly $m$-dependent trip count is absorbed by the
  block scheduler.
\item All threads of a block share one $(A_i, B_i, a_i, b_i)$ stream
  ($48$\,B per $i$): the loads are uniform across the block and hit L1/L2
  broadcast; per-transform DRAM traffic from the tables is a single
  $\sim\!0.45$\,GB stream at $N_\mathrm{side}=2048$ ($\sim$0.3\,ms at HBM
  bandwidth).  No shared memory is needed.
\item Each thread keeps $\mu_i,\mu_{i-1}$, the four accumulators
  $\mathrm{Re}/\mathrm{Im}(p_1,p_2)$, $x=\cos^2\theta$ and $\cth$ in
  registers.  The inner loop is unrolled two-fold like \texttt{ducc0}'s
  (\eqref{eq:murec} alternating the roles of $\mu_i/\mu_{i-1}$, avoiding a
  swap).
\item \textbf{Two ring pairs per thread.}  The $\mu$ recursion is a serial
  chain of two dependent FMAs per step ($\sim$8 cycles), against which each
  step offers only four independent accumulation FMAs --- one chain per
  thread cannot hide the latency even at full occupancy (measured
  $\sim$28\% of fp64 peak).  Giving each thread two \emph{adjacent} ring
  pairs (nearly equal $\theta$, hence near-identical $m_\mathrm{lim}$ and
  scale-emergence points) doubles the independent FMA streams per thread
  and amortises the shared coefficient loads, the exact GPU analogue of
  \texttt{ducc0}'s SIMD-over-rings: measured $1.45\times$ on the Legendre
  stage (66 registers, no spills).  The pair that surfaces earlier is
  advanced solo over the short gap, then both run in one fused loop.
\item Warp divergence from the scaled start-up phase is mild: threads in a
  block are neighbouring rings, whose emergence $\ell$ differs little; the
  $m_\mathrm{lim}$ skip retires whole warps at once (rings are ordered
  pole$\to$equator).
\item The flop budget at $N_\mathrm{side}=2048$, $\ell_{\max}=6143$:
  $2N_\mathrm{side}\cdot\frac14\ell_{\max}^2 \cdot 12 \approx 3.7\cdot10^{11}$
  flops (before the $\sim$30\% $m_\mathrm{lim}$ saving); A100 fp64 peak
  (non-tensor) is $9.7$\,Tflop/s.
\end{itemize}

The $(a_i,b_i)$ table is built once per plan by a tiny kernel (one thread
per $m$, sequential in $i$ --- the $\alpha$ recursion \eqref{eq:alpha} is
serial); the $(A_i,B_i)$ pre-fold \eqref{eq:prefold} runs per transform
with one thread per $(m,i)$, reconstructing
$\alpha_i = \mathrm{sign}_4(i)\sqrt{|a_i|}$ and computing $\eps$ on the fly.
\label{sec:tables}

%==============================================================================
\section{Stage 2: aliasing fold and per-ring inverse FFT}
\label{sec:fold}

Ring values follow from \eqref{eq:factorize}.  With $n \equiv n_\varphi$
pixels, $e^{im\varphi_j} = e^{im\varphi_0} e^{2\pi i\, jm/n}$ and
$e^{2\pi i\,jm/n}$ depends only on $k = m \bmod n$: every $F_m$ is
\emph{aliased} onto FFT bin $k$.  At $\ell_{\max} = 3N_\mathrm{side}-1$
this matters on \emph{all} rings, including the belt
($4N_\mathrm{side} < 2\ell_{\max}+1$).  The fold kernel computes, per ring,
\begin{equation}
G_k \;=\; \sum_{m \equiv k \ (\mathrm{mod}\ n)} t_m
     \;+\; \sum_{m \equiv -k \ (\mathrm{mod}\ n),\, m>0} t_m^{*},
\qquad t_m = F_m e^{im\varphi_0},
\end{equation}
(one thread per (ring, $m \le m_\mathrm{lim}$), \texttt{atomicAdd} into
$G$; the two conjugate contributions make $G$ Hermitian so the inverse DFT
$f_j = \sum_k G_k e^{2\pi i jk/n}$ is real).  $\mathrm{Im}\,F_0$ is
discarded, matching \texttt{ducc0}.

\textbf{Twiddle accuracy.}  $m\varphi_0$ can reach $\sim\!5000$ rad; naive
\texttt{sincos} of the rounded product would cost $\sim\!m\cdot$ulp of
phase.  Since $\varphi_0 = \pi q/d$ with integers $(q,d)$
($q{=}1, d{=}4i'$ in the caps; $q\in\{0,1\}, d{=}4N_\mathrm{side}$ in the
belt), we reduce $mq \bmod 2d$ \emph{in integer arithmetic} and call
\texttt{sincospi} --- the twiddle is then accurate to $\sim$1 ulp for every
$m$.  The same trick seeds the cap DFT rotations.

\textbf{Belt rings} all share $n = 4N_\mathrm{side}$ and are contiguous in
the RING map layout, so a single batched cuFFT inverse c2c transform of
shape $(2N_\mathrm{side}{+}1,\, 4N_\mathrm{side})$ evaluates them; a
trailing elementwise kernel writes $n\,\mathrm{Re}(\cdot)$ straight into
the map segment.

\textbf{Cap rings} have $n = 4, 8, \dots, 4(N_\mathrm{side}{-}1)$ --- all
different, awkward for cuFFT (thousands of tiny plans/launches).  Two
paths:

\emph{Small rings} ($i < 64$; negligible pixel count): direct Hermitian
inverse DFT,
\begin{equation}
f_j = G_0 + (-1)^j G_{n/2}
   + 2\sum_{k=1}^{n/2-1} \mathrm{Re}\!\left(G_k\, e^{2\pi i jk/n}\right),
\end{equation}
one thread per pixel, rotating $e^{2\pi i jk/n}$ incrementally and
re-seeding it exactly (integer-reduced \texttt{sincospi}) every 128 steps,
which bounds the phase drift at $\sim\!10^{-14}$.  (Applied to \emph{all}
cap rings this kernel is $O(N_\mathrm{side}^3)$ and was $\sim$60\% of the
transform at $N_\mathrm{side}=2048$; it survives only for the tiny rings
where Bluestein padding would be wasteful.)

\emph{All other cap rings}: Bluestein's algorithm through batched cuFFT.
With chirps $w_t = e^{i\pi t^2/n}$ (computed exactly: $t^2 \bmod 2n$ in
integers, then \texttt{sincospi}),
\begin{equation}
f_j \;=\; w_j \sum_{k} \left(G_k w_k\right)\, \bar w_{j-k},
\end{equation}
a linear convolution evaluated as a circular one of power-of-two length
$M \ge 2n-1$.  Rings are grouped into size classes
$i \in [2^j, 2^{j+1})$ sharing $M = 2^{j+4}$, so \emph{one} batched
cuFFT pair (forward + inverse) per class handles a whole hemisphere pair
population; the FFTs of the chirp kernels $\hat b$ are precomputed at plan
build (715\,MB at $N_\mathrm{side}=2048$).  This replaces the
$O(N_\mathrm{side}^3)$ cap cost by $O(N_\mathrm{side}^2 \log
N_\mathrm{side})$ and cut the $N_\mathrm{side}=2048$ transform from
179\,ms to 112\,ms.

%==============================================================================
\section{Accuracy}
\label{sec:accuracy}

Correctness is defined as agreement with \texttt{ducc0}
(\texttt{synthesis}, RING geometry from \texttt{Healpix\_Base.sht\_info()}).
Because Almond implements the \emph{same} recursion, coefficients, scaling
thresholds and $m_\mathrm{lim}$ rule, agreement is at the level of
floating-point noncommutativity, not method error.  Gates in the test
suite (random Gaussian $a_{\ell m}$, $\ell_{\max}=3N_\mathrm{side}-1$):
\begin{center}
\begin{tabular}{lcc}
\toprule
check & criterion & measured \\
\midrule
NumPy reference vs \texttt{ducc0}, $N_\mathrm{side}\!\le\!32$ &
  $<10^{-12}$ & $\sim10^{-14}$ \\
GPU vs \texttt{ducc0}, $N_\mathrm{side} = 8\ldots1024$ &
  $<10^{-10}$ & \textit{(see \S\ref{sec:results})} \\
call-to-call determinism (atomics reorder) & $\le10^{-13}$ & pass \\
\bottomrule
\end{tabular}
\end{center}
The transform is deterministic up to the order of the \texttt{atomicAdd}s
in the fold stage, whose additions involve $\lesssim m_{\max}/n$ terms per
bin; the observed call-to-call scatter is $\lesssim10^{-15}$ relative.

A pure-NumPy reference implementation (\texttt{almond/reference.py}), ported
line-by-line from \texttt{ducc0} and validated against it first, serves as
the executable specification: every GPU kernel corresponds one-to-one to a
block of that file, which made the CUDA port correct on first run.

%==============================================================================
\section{Memory budget}
\label{sec:memory}

Persistent device buffers of a plan at $(N_\mathrm{side}, \ell_{\max} =
3N_\mathrm{side}-1)$, in double precision:

\begin{center}
\begin{tabular}{lll}
\toprule
buffer & size & $N_\mathrm{side}=2048$ \\
\midrule
coefficient table $(a,b)$ & $16\,\mathrm{B}\cdot\ell_{\max}^2/4$ & 151 MB \\
pre-folded $(A,B)$        & $32\,\mathrm{B}\cdot\ell_{\max}^2/4$ & 302 MB \\
phase $F_m(\mathrm{ring})$ & $16\,\mathrm{B}\,(m_{\max}{+}1)N_\mathrm{ring}$ & 805 MB \\
FFT bins $G$              & $16\,\mathrm{B}\,N_\mathrm{pix}$ & 805 MB \\
output map                & $8\,\mathrm{B}\,N_\mathrm{pix}$ & 403 MB \\
Bluestein $\hat b$ tables & $\sim 16\,\mathrm{B}\cdot 43 N_\mathrm{side}^2$ & 715 MB \\
\midrule
total persistent          & & $\sim$3.2 GB \\
transients (belt + Bluestein cuFFT) & & $\sim$1.8 GB \\
\bottomrule
\end{tabular}
\end{center}
plus the input $a_{\ell m}$ (302 MB) if resident.  Comfortable on a 40 GB
A100 up to $N_\mathrm{side} = 4096$; the phase and $G$ buffers could be
eliminated altogether by fusing the fold into the Legendre kernel
(\S\ref{sec:future}).

%==============================================================================
\section{Benchmarks}
\label{sec:results}

Measured on a dedicated Perlmutter GPU node (A100-SXM 40\,GB; AMD EPYC
7763, 64 physical cores), SLURM job 55427184.  Methodology: accuracy gate
vs \texttt{ducc0} first; then median of $\ge$5--10 warm calls with
\texttt{deviceSynchronize} around each; \texttt{ducc0} timed on the same
node at 1/16/32/64 threads (best time reported); Almond timed both
device-resident (the QML use case --- the CG solver keeps data on the GPU)
and including host$\to$device$\to$host copies of $a_{\ell m}$ and the map.
Random Gaussian $a_{\ell m}$, $\ell_{\max} = 3N_\mathrm{side}-1$, spin 0,
single transform ($B=1$).

\begin{center}
\begin{tabular}{rrcrrrrrr}
\toprule
$N_\mathrm{side}$ & $\ell_{\max}$ & acc.\ vs ducc & Alm dev.\ [s] & Alm host [s] & ducc 1t [s] & ducc best [s] & speedup & peak GB \\
\midrule
128 & 383 & $2\times10^{-13}$ & 0.0003 & 0.0006 & 0.009 & 0.001 (32t) & 2.6$\times$ & 0.02 \\
256 & 767 & $4\times10^{-13}$ & 0.0005 & 0.0014 & 0.053 & 0.002 (64t) & 4.6$\times$ & 0.08 \\
512 & 1535 & $8\times10^{-13}$ & 0.0023 & 0.0060 & 0.353 & 0.013 (64t) & 5.7$\times$ & 0.31 \\
1024 & 3071 & $3\times10^{-12}$ & 0.0133 & 0.0291 & 2.554 & 0.081 (64t) & 6.1$\times$ & 1.24 \\
2048 & 6143 & $6\times10^{-12}$ & 0.0770 & 0.1558 & 19.110 & 0.479 (64t) & 6.2$\times$ & 4.97 \\
\bottomrule
\end{tabular}

\end{center}

\noindent Reading of the table:
\begin{itemize}
\item \textbf{Almond beats best-threaded (64-core) ducc0 at every resolution},
  by $6.2\times$ at $N_\mathrm{side}=2048$ (77\,ms vs 479\,ms;
  $250\times$ vs single-thread ducc0).  Against the QML-relevant
  device-resident workload there is no CPU/GPU crossover to wait for ---
  the GPU wins already at $N_\mathrm{side}=128$.
\item Including PCIe copies the win shrinks to $1.2$--$3.1\times$:
  transfers cost roughly as much as the transform itself at large
  $N_\mathrm{side}$ (455\,MB round trip at $N_\mathrm{side}=2048$).
  Batched/pipelined transfers or keeping data resident (SiMaster's mode)
  are essential.
\item Whole-map throughput at $N_\mathrm{side}=2048$:
  $50.3\,\mathrm{Mpix} / 77\,\mathrm{ms} = 653$\,Mpix/s.  Stage split:
  Legendre 65\,ms ($\sim$4\,Tflop/s, 40\% of A100 fp64 peak), Bluestein
  caps 7.6\,ms, fold 4.9\,ms, belt cuFFT 3.3\,ms, pre-fold 0.6\,ms.
\item Accuracy vs ducc0 degrades slowly with $\ell_{\max}$
  ($1.9\times10^{-13} \to 5.6\times10^{-12}$), consistent with
  floating-point noncommutativity of the fold atomics and the different
  twiddle paths, far below the $10^{-8}$ science gate.
\item Peak device memory (plan + transients): 5.0\,GB at
  $N_\mathrm{side}=2048$, dominated by the phase/$G$ buffers and Bluestein
  tables (\S\ref{sec:memory}); $N_\mathrm{side}=4096$ fits a 40\,GB A100.
\end{itemize}

\subsection{High-throughput batched comparison}
\label{sec:batched}

QML applies the SHT to many right-hand sides, and \texttt{ducc0}'s
strongest many-transform configuration is not internal threading of one
transform but its \texttt{ntrans} batch mode ($a_{\ell m}$ of shape
$(B, 1, n_{\rm alm})$, one call, threads across the batch).  On few-core
runs this mode is dramatically better per column than intra-transform
threading (login-node check: $426 \to 41.6$\,ms/col from 1 to 16 threads at
$N_\mathrm{side}=512$).  The honest throughput comparison is therefore
per-column time at large $B$ against \texttt{ducc0}-ntrans at 64 threads
(same dedicated node; Almond runs \texttt{synthesis\_device\_batch}):

\begin{center}
\begin{tabular}{rrcrrrr}
\toprule
$N_\mathrm{side}$ & $B$ & acc. & Alm [ms/col] & Alm+copies & ducc0 64t [ms/col] & speedup \\
\midrule
256 & 128 & $6\times10^{-13}$ & 0.39 & 1.55 & 2.11 & 5.4$\times$ \\
512 & 64 & $1\times10^{-12}$ & 1.98 & 6.52 & 12.01 & 6.1$\times$ \\
1024 & 64 & $4\times10^{-12}$ & 11.26 & 31.12 & 68.88 & 6.1$\times$ \\
2048 & 16 & $7\times10^{-12}$ & 76.96 & 150.27 & 460.96 & 6.0$\times$ \\
\bottomrule
\end{tabular}

\end{center}

Two findings.  First, on a full 64-core node \texttt{ducc0}'s batch mode is
barely faster per column than its (well-implemented) intra-transform
threading (461 vs 479\,ms/col at $N_\mathrm{side}=2048$): 64 concurrent
CPU transforms saturate host memory bandwidth long before they saturate
64 cores.  Almond's advantage in the throughput regime is therefore the same
$\sim\!6\times$ as for single transforms, uniformly over
$N_\mathrm{side} = 256$--$2048$ (1.4--3.1$\times$ if every batch crosses
PCIe both ways).

Second, a negative result on the GPU side: a chunked Legendre kernel that
shares one $\mu$ recursion among $C$ columns (4$C$ accumulators per
thread) \emph{loses} to simply looping the single-transform pipeline
(11.7\,ms/col sequential vs 14.4--16.7\,ms/col for $C = 2$--$8$ at
$N_\mathrm{side}=1024$): the accumulator registers destroy occupancy (150
registers at $C=8$) and each column needs its own $(A_i,B_i)$ loads,
doubling the relative L1 pressure, while the amortised recursion was only
ever 4 of 12 flops per step.  The two-ring-pair kernel of
\S\ref{sec:kernel} already owns the right trade: its two chains share one
coefficient stream.  Ring pairs beat column chunks on the GPU for the same
reason SIMD-over-rings beats SIMD-over-columns on the CPU.  The batched
API therefore defaults to the looped single pipeline (asynchronous, no
per-column synchronisation); the chunked kernel is retained behind
\texttt{chunk>1} for future spin-2/adjoint experiments.

%==============================================================================
\section{Design decisions and departures from ducc0}
\label{sec:departures}

\texttt{ducc0} is aggressively CPU-shaped; porting it verbatim would waste
the GPU.  What was kept, and what was changed:

\begin{itemize}
\item \textbf{Kept: the mathematics.}  Recursion \eqref{eq:murec},
  pre-fold \eqref{eq:prefold}, scaling thresholds, $m_\mathrm{lim}$ ---
  bit-compatible behaviour makes validation trivial and inherits 15 years
  of numerical hardening.
\item \textbf{Changed: parallel decomposition.}  CPU: SIMD over
  $\sim$8--64 rings within one thread, OpenMP over $m$, ring-block cache
  tiling (\texttt{lstep}) to keep $a_{\ell m}$ resident in L2.  GPU: one
  thread per (ring pair, $m$), $2.5\cdot10^7$-way parallelism, no tiling
  --- the coefficient streams are block-uniform loads that L1/L2 serve at
  register speed, and the $\sim$150 MB tables are read once per transform
  from HBM.
\item \textbf{Changed: coefficient handling.}  \texttt{ducc0} recomputes
  the per-$m$ coefficient vectors on each core inside the $m$ loop (cheap
  on CPU, cache-resident).  We precompute the full triangular $(a,b)$
  table once per plan (151 MB at $N_\mathrm{side}{=}2048$): GPU threads
  cannot afford a serial per-$m$ preparation, and memory is cheaper than
  divergence.  $\alpha_i$ is \emph{not} stored --- it is reconstructed as
  $\mathrm{sign}_4(i)\sqrt{|a_i|}$, saving 75 MB.
\item \textbf{Changed: FFT stage.}  \texttt{ducc0} runs pocketfft
  ring-by-ring inside the OpenMP loop with FFTPACK-style half-complex
  buffers.  We split by ring population: one batched cuFFT for the
  $2N_\mathrm{side}{+}1$ equal-length belt rings, a direct Hermitian DFT
  kernel for the caps, and exact-rational \texttt{sincospi} twiddles
  instead of a cached \texttt{MultiExp} table.
\item \textbf{Changed: precision management is per-thread scalar.}  The
  SIMD-lane \texttt{corfac} machinery collapses to a simple three-phase
  scalar loop (skip / unscale / plain FMA), which is both simpler and
  divergence-friendly.
\end{itemize}

%==============================================================================
\section{Adjoint synthesis and spin-2 (v0.2)}
\label{sec:v02}

\textbf{Adjoint} ($Y^\dagger$: map$\to$alm, ducc0's convention --- the
transpose under the alm inner product that counts $m>0$ twice): per ring
$\hat G_k = \sum_j f_j e^{-2\pi i jk/n}$ (cuFFT rfft on the belt,
sign-flipped Bluestein on the caps), unfold
$F'_m = e^{-im\varphi_0}\hat G_{m \bmod n}$ (with $\mathrm{Im}\,F'_0 \to
0$), then the Legendre adjoint
$A'_i = \sum_\mathrm{pairs} \mu_i\,(F'_N{+}F'_S)$,
$B'_i = \sum_\mathrm{pairs} \mu_i \cth (F'_N{-}F'_S)$ and an
$\alpha/\eps$-weighted gather to $a'_{\ell m}$.  The reduction kernel gives
each $m$ one thread block owning \emph{all} ring pairs (chunk loop), so
global accumulation needs no atomics: per step, contributions are
warp-shuffle-reduced (4 lanes-worth of chains per lane) and staged in
per-warp shared tiles flushed every 128 steps.  Validated against
\texttt{ducc0.adjoint\_synthesis} to $<10^{-10}$ (nside 8--1024) and
against Almond's own synthesis via the transpose identity to $10^{-12}$.
Cost: $\sim$1.5--1.9$\times$ the synthesis (147 vs 77\,ms at
$N_\mathrm{side}{=}2048$) --- the shuffle-reduction overhead.

\textbf{Spin-2}: the native \texttt{ducc0} spin recursion (its
spin-2-via-spin-0 shortcut with $1/\sin^2\theta$ ring weights is inexact
near the poles; ducc itself reverts to the native path for
$\sin\theta<0.01$).  Two scaled Wigner-$d$ chains
$(\lambda^\pm \sim d^\ell_{m,\mp 2})$ advance with a $\Delta\ell=1$
three-term recursion in $\cth$; eight accumulators build the $Q\pm iU$
combinations with parity-interleaved $E/B$ coupling; initialisation uses
half-angle power prefactors $\mathrm{prefac}(m)\,c^{a}s^{b}$ carried in the
$2^{800}$-scaled representation.  Stage 2 (fold/FFT) is reused verbatim
with the two components riding the batch-chunk axis.  Synthesis and
adjoint validated against ducc0 to $<10^{-10}$ (nside 8--512).

Two porting bugs are worth recording.  (i) prefac$(m)$ underflowed to
exactly zero for $m \gtrsim 514$: the factorial-table \emph{ratios} need
ducc's two-sided normalisation --- one-sided clamping assumes a direction
values do not actually keep.  Same lesson as the \texttt{mypow} incident:
port the invariant band, not the assumed monotonicity.  (ii) the spin
adjoint requires \emph{conjugated} coefficients on the $\pm i$ cross terms
(real-linear transpose: $y \mathrel{+}= c\,x \Rightarrow x' \mathrel{+}=
\bar c\,y$).  Both were found by mode-range bisection against ducc0.

\textbf{Measured} (dedicated A100-SXM + 64-core EPYC, job 55449334;
device-resident, median of warm calls; ducc0 at 64 threads; accuracy gate
passed in all cells):

\begin{center}
\begin{tabular}{rcrrrr}
\toprule
$N_\mathrm{side}$ & acc.\ (spin 2) & \multicolumn{4}{c}{Almond [ms] (speedup vs ducc0 64t)} \\
 & & synth$_0$ & adj$_0$ & synth$_2$ & adj$_2$ \\
\midrule
512 & $3\times10^{-12}$ & 2.0 (6.3$\times$) & 3.7 (3.3$\times$) & 9.0 (3.3$\times$) & 14.5 (1.8$\times$) \\
1024 & $6\times10^{-12}$ & 12.1 (6.5$\times$) & 25.0 (3.2$\times$) & 49.4 (3.0$\times$) & 80.5 (1.7$\times$) \\
2048 & $2\times10^{-11}$ & 71.9 (6.5$\times$) & 139.7 (3.2$\times$) & 363.5 (2.6$\times$) & 587.4 (1.6$\times$) \\
\bottomrule
\end{tabular}

\end{center}

\noindent The adjoint speedups ($\sim$3.2$\times$ spin-0, $\sim$1.7$\times$
spin-2) trail the synthesis ($6.5\times$ / $\sim$3$\times$) because ducc0's
adjoint costs the same as its synthesis while Almond's reduction kernels
pay the shuffle overhead --- the known tuning target.

\textbf{SiMaster integration}: \texttt{almond.simaster.AlmondRealSHT} is a
validated drop-in for \texttt{simaster.sht.RealSHT} (real coefficient
basis, observed-pixel subsetting, batched columns; numpy or cupy in/out),
tested against it for spin 0 and 2 on a cut sky to $10^{-10}$ together
with the real-basis transpose identity.

\section{v0.3: fused fold, tuning, SiMaster integration}
\label{sec:v03}

The fold/unfold stages are now \emph{fused into the Legendre kernels}: the
forward extract scatters the twiddled $F_m$ straight onto the ring FFT
bins (\texttt{fold\_scatter}, atomicAdd), and the adjoint init gathers
$F'_m$ from the half-spectrum $\hat G$ on the fly.  The phase array no
longer exists in any single-transform path ($-805$\,MB spin-0 /
$-1.6$\,GB spin-2 at $N_\mathrm{side}=2048$, plus one full memory pass in
each direction); runtimes are unchanged within noise, as expected --- this
was a memory optimisation.  The spin-2 adjoint reduction now runs four
chains per lane (664$\to$596\,ms at $N_\mathrm{side}=2048$; 168 registers,
occupancy-limited).

SiMaster integration is complete short of the physics reruns:
\texttt{almond} is an editable package and
\texttt{CovModel}/\texttt{QMLWorkspace} accept \texttt{backend='almond'}
(same \texttt{pure\_callback} plumbing as \texttt{'ducc'}).  The
covariance operator $C\!\cdot\!x$ agrees between the two backends to
$10^{-10}$ on a cut sky with mixed spin-0/spin-2 fields, and SiMaster's
own test suite passes unchanged.

\section{Roadmap}
\label{sec:future}

\begin{enumerate}
\item \textbf{val1/val2 QML validation} with \texttt{backend='almond'}
  (deliberately left to a SiMaster-focused follow-up).
\item \textbf{Spin-2 adjoint reduction restructure}: SADJ\_PPT=4 is
  register-bound; a block-level two-stage reduction could recover the
  remaining $\sim$1.5$\times$ to parity with the forward kernel.
\item \textbf{Batched accumulator chunking} (one recursion feeding $C$
  columns) \emph{measured to lose} to looping the single-transform
  pipeline (registers + L1 pressure; see \S\ref{sec:batched}).
\end{enumerate}

\end{document}
